\TopicDivider{01}{Prior, likelihood, posterior}{The update is a rule for combining what was believed with what was observed.}

\begin{frame}{Learning goals}
  \begin{LearningGoals}{By the end of this topic, learners can}
    \begin{itemize}
      \item distinguish a prior, likelihood and posterior in a one-parameter model;
      \item state Bayes' rule without treating the denominator as decoration; and
      \item explain why a posterior is a distribution rather than a single answer.
    \end{itemize}
  \end{LearningGoals}
\end{frame}

\begin{frame}{The update has three interpretable parts}
  \TwoPane{
    \begin{DefinitionCard}{Prior}
      A distribution encoding uncertainty about a parameter before the current evidence is observed.
    \end{DefinitionCard}
    \vspace{.22cm}
    \begin{DefinitionCard}{Likelihood}
      A model of how plausible the observed evidence is under each possible parameter value.
    \end{DefinitionCard}
  }{
    \begin{DefinitionCard}{Posterior}
      The updated distribution after the prior and likelihood have been combined and normalised.
    \end{DefinitionCard}
    \vspace{.22cm}
    \begin{KeyTakeaway}{Keep the objects separate}
      A likelihood is not a probability distribution over the parameter until it is combined with a prior and normalised.
    \end{KeyTakeaway}
  }
\end{frame}

\begin{frame}{Bayes' rule makes the assumption visible}
  \BayesEquation{p(\theta\mid y)=\frac{p(y\mid\theta)\,p(\theta)}{p(y)}}
  \begin{ProofSketch}{Read the denominator as a model check}
    The evidence \(p(y)\) ensures the posterior integrates to one. It also rewards models that assign reasonable probability to the data actually observed.
  \end{ProofSketch}
\end{frame}

\begin{frame}{TeXScience plot: an update shifts mass, not certainty}
  \begin{ScientificFigure}{Illustrative one-dimensional densities}
    \centering
    \begin{tikzpicture}
      \begin{axis}[bayes/science plot,xlabel={parameter value},ylabel={relative density},xmin=0,xmax=10,ymin=0,ymax=1.08,legend pos=north west]
        \addplot[color=BayesMuted,domain=0:10,samples=100] {exp(-.13*(x-3)^2)};
        \addlegendentry{prior}
        \addplot[color=BayesTeal,domain=0:10,samples=100] {exp(-.34*(x-6.1)^2)};
        \addlegendentry{posterior}
      \end{axis}
    \end{tikzpicture}
    \BayesCaption{PGFPlots source stays editable in version control; the distribution shape is illustrative.}
  \end{ScientificFigure}
\end{frame}
